\documentclass[11pt]{article}

\usepackage[margin=2.2cm]{geometry}
\usepackage{xcolor}
\usepackage{forest}
\usepackage{tikz}
\usetikzlibrary{arrows.meta,positioning,decorations.pathreplacing,calc}
\usepackage{amsmath,amssymb,mathtools}
\usepackage{booktabs}
\usepackage{array}
\usepackage{enumitem}
\usepackage{titlesec}
\usepackage{float}
\usepackage[colorlinks=true,linkcolor=teal!70!black,citecolor=teal,urlcolor=teal]{hyperref}
\usepackage{parskip}
\usepackage{cleveref}

\setlength{\parskip}{0.5\baselineskip}
\titlespacing*{\section}{0pt}{1.5em}{0.6em}
\titlespacing*{\subsection}{0pt}{1.2em}{0.4em}

% ---- automatic "Part N --- Title" section numbering, so \section + \label
% gives both the on-page heading and a working \ref/\cref everywhere else ----
\renewcommand{\thesection}{Part \arabic{section}}
\titleformat{\section}{\normalfont\Large\bfseries}{\thesection}{0.6em}{---\ }

% ---- colors reused across every figure, so the same idea always wears the same color ----
\definecolor{lossless}{HTML}{1f7a6c}   % teal family
\definecolor{lossy}{HTML}{c2622d}      % orange family
\definecolor{prefixcol}{HTML}{6a4c93}  % violet
\definecolor{nonprefixcol}{HTML}{1f6fb2}% blue
\definecolor{dictcol}{HTML}{8a5a2b}    % brown
\definecolor{bestcol}{HTML}{2e8b3d}    % green, "best in class" marker
\definecolor{warncol}{HTML}{b3261e}    % red, "error / impossible" marker
\definecolor{transformcol}{HTML}{2f7d8f}% steel teal, transform-based family
\definecolor{contextcol}{HTML}{a13d63} % magenta, context-modeling family
\definecolor{mmcol}{HTML}{4f6d2f}      % olive, lossless multimedia family
\definecolor{groupcol}{HTML}{1f6fb2}   % blue, used for the group axioms box

\newcommand{\figdrawnote}[1]{\par\noindent\textbf{\textcolor{gray!60!black}{[Draw on the board here:}} \textit{#1}\textbf{\textcolor{gray!60!black}{]}}\par}

% ---- a colored callout box for direct-to-camera, pause-and-think moments ----
\newcommand{\engage}[3]{%
\begin{center}
\begin{tikzpicture}
\node[draw=#1, thick, rounded corners, fill=#1!6, align=left,
      font=\sffamily\small, text width=0.92\linewidth, inner sep=9pt]
      {\textbf{\textcolor{#1}{#2}}\; #3};
\end{tikzpicture}
\end{center}
}

% ---- a light gray box for verbatim "say this to camera" narration ----
\newcommand{\say}[1]{%
\begin{center}
\begin{tikzpicture}
\node[draw=gray!50, thick, rounded corners, fill=gray!5, align=left,
      font=\itshape, text width=0.92\linewidth, inner sep=9pt]
      {#1};
\end{tikzpicture}
\end{center}
}

\title{\Large\bfseries What is a Group?\\[2pt]
\normalsize\normalfont\itshape The first stop on our tour of abstract algebra}
\author{Dr. Faisal Aslam}
\date{}

\begin{document}
\maketitle
\vspace{-2em}

\say{Hi everyone, welcome back. Today we're starting a brand new topic: \textbf{abstract algebra} --- and our very first stop is an object called a \textbf{group}. Groups are used in cryptography, computer graphics, and coding theory. However, our main interest is in cryptography.

So here's the plan of this video. First, I'll define the group and then provide many examples to make sure everything is clear. Let's get into it.}

\section{Groups}
\label{sec:L1:groups}

\say{Alright, let's start with the definition.}

A set $G$ together with a binary operation $*$ is called a \textbf{group}, written $(G, *)$. Here $*$ can be \emph{any} binary operation --- ordinary multiplication, addition, XOR, matrix multiplication, function composition, or even some strange custom-made rule you invented yourself, as long as it takes two elements of $G$ and produces a single output. Likewise the set $G$ can contain anything: numbers, vectors, matrices, polynomials, symmetries of a shape, or any custom-built objects at all. What makes $(G,*)$ a \emph{group} is not what the elements are, but whether the pair $(G,*)$ obeys the following four rules, called \textbf{axioms}.

A group must satisfy all four of the following axioms:
\begin{enumerate}[leftmargin=*, label=(G\arabic*)]
  \item \textbf{Closure:} for all $a, b \in G$, $a * b \in G$.
  \item \textbf{Associativity:} for all $a,b,c \in G$, $(a*b)*c = a*(b*c)$.
  \item \textbf{Identity:} there exists an element $e \in G$ such that $a * e = e * a = a$ for all $a \in G$.
  \item \textbf{Inverses:} for every $a \in G$, there exists $a^{-1} \in G$ such that $a * a^{-1} = a^{-1} * a = e$.
\end{enumerate}

Note that the identity element is common to \emph{every} element of the group, whereas each element of the group may have its own, possibly different, inverse. If, in addition to the four axioms above, the group's binary operation is also \textbf{commutative} --- that is, $a * b = b * a$ for all $a, b \in G$, so the order in which you combine two elements doesn't matter --- then the group is called an \textbf{abelian group}, named after the Norwegian mathematician Niels Henrik Abel.

\subsection{Examples}
Let's cement our understanding of groups with a range of examples --- including a few near-misses, because seeing \emph{why} something fails an axiom teaches you the axioms far better than only seeing successes ever could.

\begin{itemize}[leftmargin=*]
  \item $(\mathbb{Z}, +)$ --- the integers under addition. Identity is $0$; the inverse of $a$ is $-a$. Adding two integers always gives an integer, addition of integers is associative, and this is an abelian group. \say{This is the group hiding behind our clock warm-up, minus the wrap-around.}

  \item $(\mathbb{R} \setminus \{0\}, \times)$ --- nonzero real numbers under multiplication. Identity is $1$; the inverse of $a$ is $1/a$. Also abelian. (We had to remove $0$ because $0$ has no multiplicative inverse --- there's no real number you can multiply $0$ by to get back to $1$.)

  \item $(\mathbb{Z}, \times)$ --- integers under multiplication --- is \textbf{not} a group: $2$ has no inverse in $\mathbb{Z}$ (we'd need $1/2$, which isn't an integer). This fails axiom (G4). \say{So close, and yet so far --- closure, associativity, and identity ($1$) all hold. It's a single missing axiom that keeps this out of the club.}

  \item \textbf{$(\mathbb{Z}_{12}, + (\mod 12))$, the integers mod 12 under addition.} The group contains integers $\{0,\dots,11\}$. The operation of addition is followed by a mod 12 operation. So $10+4 = 14 \equiv 2 \pmod{12}$, matching our warm-up exactly. The identity is $0$, and the inverse of $a$ is $12-a$ (check: $10 + 2 = 12 \to 0$). This is an abelian group, and more generally $(\mathbb{Z}_n,+)$ is a group for \emph{every} whole number $n \ge 1$ --- these are called \textbf{cyclic groups}, and they quietly run underneath almost every hashing and error-detection scheme you've ever used.

  \item \textbf{$(\mathbb{Z}_{p} \setminus \{0\}, \times \pmod p)$, where $p$ is a prime number.} This is a group under multiplication. These groups have lots of applications in cryptography.

  \item \textbf{$3\times 3$ real matrices under addition, $(\mathbb{R}^{3\times 3}, +)$, is a group.} Adding two $3\times 3$ matrices gives another $3\times 3$ matrix (closure); matrix addition is associative; the identity is the all-zeros matrix $\mathbf{0}$; the inverse of $A$ is $-A$ (negate every entry). It's even abelian, since $A+B=B+A$ entrywise. Nothing about this is different in spirit from $(\mathbb{Z},+)$ --- just $9$ numbers moving in lockstep instead of $1$.

  \item \textbf{The same set under multiplication is \emph{not} a group.} Try $(M_{3\times 3}(\mathbb{R}), \times)$ (ordinary matrix multiplication): closure holds (multiplying two $3\times3$ matrices gives another $3\times3$ matrix) and multiplication is associative, and there's an identity, $I_3$. But axiom (G4) fails: a \textbf{singular} matrix (one with $\det = 0$, e.g.\ the all-zeros matrix, or $\begin{pmatrix}1&0&0\\0&1&0\\0&0&0\end{pmatrix}$) has no inverse matrix at all --- there is no $B$ with $AB=BA=I_3$. So the full set of $3\times3$ matrices under multiplication is not a group, for exactly the same structural reason $(\mathbb{Z},\times)$ wasn't: some elements simply lack an inverse.

  \item \textbf{Invertible matrices under multiplication, on the other hand, \emph{do} form a group.} Restrict to only the invertible $3\times 3$ matrices, $GL_3(\mathbb{R}) = \{A \in M_{3\times3}(\mathbb{R}) : \det(A)\ne 0\}$. Now every element has an inverse \emph{by construction} (that's what ``invertible'' means), the product of two invertible matrices is again invertible (closure restored), and $I_3$ is still the identity. This is called the \textbf{general linear group}, and it's the standard fix for the multiplication problem above: throw away the elements that were breaking (G4), and check that what's left is still closed. \say{This is also our first example of a \textbf{non-abelian} group: in general $AB \neq BA$ for matrices, so order really does matter here --- swap the order of two moves and you can land somewhere completely different. Video game engines and robotics software run into this constantly: rotating an object and then moving it is \emph{not} the same as moving it and then rotating it.}

  \item \textbf{The symmetries of a square, formally: the dihedral group $D_4$.} \figdrawnote{the square from the warm-up again, now also showing 4 flip lines: horizontal, vertical, and the two diagonals.} Earlier we only rotated the square. If we also allow \emph{flipping} it over (a reflection), we get $8$ symmetries in total: $4$ rotations ($0^\circ,90^\circ,180^\circ,270^\circ$) and $4$ reflections. Combining ``do symmetry $a$, then symmetry $b$'' is the group operation, the identity is ``do nothing,'' and every symmetry can be undone by another symmetry in the set (rotate $90^\circ$ forward, undo by rotating $270^\circ$ forward; flip along an axis, undo by flipping along the \emph{same} axis again). This group is called $D_4$, and crucially, it is \textbf{not abelian}: rotate $90^\circ$ then flip horizontally, versus flip horizontally then rotate $90^\circ$, land the square in two different final positions. Try it with a piece of paper --- it's the fastest way to feel non-commutativity in your hands.

  \item \textbf{Degree-$<n$ polynomials under addition form a group; under multiplication they do not.} Let $P_n = \{a_0+a_1x+\cdots+a_{n-1}x^{n-1} : a_i \in \mathbb{R}\}$ (real polynomials of degree strictly less than $n$). Under addition, $(P_n,+)$ is a group exactly like matrices under addition above: closed (adding two such polynomials keeps the degree below $n$), associative, identity is the zero polynomial, and the inverse of $p(x)$ is $-p(x)$. Under multiplication, it immediately fails on \emph{two} separate grounds: closure fails (multiplying two degree-$(n-1)$ polynomials can produce degree up to $2n-2$, which is no longer in $P_n$ at all), and even ignoring that, most individual polynomials have no multiplicative inverse within polynomials (e.g., there is no polynomial $q(x)$ with $x\cdot q(x) = 1$ for every $x$). \say{Hold on to this particular example --- one operation behaving like a group and the other one refusing to is exactly the gap that motivates our \emph{next} lecture, on a structure called a \textbf{ring}, where we finally allow two operations to coexist without demanding that both of them be full groups.}
\end{itemize}


\subsection{How big is a group? A quick word on order}

\say{One more piece of vocabulary before we move to applications. Some of our groups, like $\mathbb{Z}$ or $\mathbb{R}\setminus\{0\}$, have infinitely many elements. Others, like our modolus group $\mathbb{Z}_{12}$ or the square's symmetry group $D_4$, have only finitely many. The number of elements in a group $G$ is called its \textbf{order}, written $|G|$. So $|\mathbb{Z}_{12}| = 12$ and $|D_4| = 8$. This single number turns out to control a surprising amount of a group's behavior, and it will come back again and again once we start talking about subgroups.}

\subsection{Why Groups Show Up Everywhere in Cryptography}

Many classical cryptographic schemes --- Diffie--Hellman key exchange, ElGamal encryption, and elliptic-curve cryptography --- are built directly on top of group structure. Concretely, Diffie--Hellman works inside a group like $(\mathbb{Z}_p^*, \times)$, the nonzero numbers modulo a large prime $p$ under multiplication. Two people can each pick a secret exponent, publicly exchange $g^a \bmod p$ and $g^b \bmod p$, and each computes the same shared secret $g^{ab} \bmod p$ --- while an eavesdropper who only sees $g$, $g^a$, and $g^b$ is stuck with what's called the \textbf{discrete logarithm problem}: recovering $a$ from $g^a$ inside this group is believed to be computationally very hard. Elliptic-curve cryptography (the scheme protecting most of the internet's HTTPS traffic today) replaces $\mathbb{Z}_p^*$ with a cleverer group built from points on a curve, but the underlying idea --- a hard problem living inside a group --- is identical. Lattice-based cryptography, one of the leading candidates for staying secure against future quantum computers, instead builds on richer structures (rings and modules, coming in our next lecture), but a group is always the first layer underneath.

\say{That's it for today. If any of the four axioms still feels shaky, rewind and try the clock and the square examples again with a pen in hand --- abstract algebra rewards you for getting your hands dirty with small, concrete examples before trusting the general definition. See you in the next one.}

\end{document}
